\documentclass[11pt]{amsart}
\usepackage[T1]{fontenc}
\usepackage{lmodern,microtype,amsmath,amssymb,amsthm,mathtools,enumitem,booktabs}
\usepackage[margin=1.1in]{geometry}
\usepackage[hidelinks]{hyperref}
\usepackage{xurl}

\newtheorem{theorem}{Theorem}[section]
\newtheorem{lemma}[theorem]{Lemma}
\newtheorem{corollary}[theorem]{Corollary}
\newtheorem{proposition}[theorem]{Proposition}
\theoremstyle{remark}
\newtheorem{remark}[theorem]{Remark}
\numberwithin{equation}{section}

\newcommand{\E}{\mathbb E}
\newcommand{\Pp}{\mathbb P}
\newcommand{\Var}{\operatorname{Var}}
\newcommand{\Cov}{\operatorname{Cov}}
\newcommand{\dd}{\,\mathrm d}

\hypersetup{pdftitle={Correlation of leaf heights at mesoscopic separation in the critical beta-splitting tree},pdfauthor={Yitzchak Shmalo}}

\title[Leaf heights at mesoscopic separation]{Correlation of leaf heights at mesoscopic separation in the critical beta-splitting tree}
\author{Yitzchak Shmalo}
\address{Einstein Institute of Mathematics, The Hebrew University of Jerusalem, Jerusalem, Israel}
\date{15 October 2026}
\subjclass[2020]{60J80, 60C05, 05C05}
\keywords{critical beta-splitting tree, leaf height, fragmentation, correlation, martingale}

\begin{document}
\begin{abstract}
In the critical beta-splitting random tree on $n$ leaves, in its interval-splitting form with the interval closed into a cycle, let $D^{(1)}_n$ and $D^{(2)}_n$ be the continuous heights of a uniformly random leaf and of the leaf $r$ places further round the cycle. We show that, uniformly in $1\le r\le n/2$,
\[
\Cov\bigl(D^{(1)}_n,D^{(2)}_n\bigr)=\sigma^2\log(n/r)+O\bigl(\sqrt{\log n}\bigr),\qquad \Var(D_n)=\sigma^2\log n+O\bigl(\sqrt{\log n}\bigr),
\]
with $\sigma^2=2\zeta(3)/\zeta(2)^3$. Hence the correlation coefficient equals $1-\log r/\log n+O((\log n)^{-1/2})$, and it tends to $1-\alpha$ whenever $\log r_n/\log n\to\alpha\in[0,1]$. This answers Open Problem~1 of Aldous and Janson. The proof follows a positional potential of each leaf, which decreases at rate $\zeta(2)$ up to errors that are square integrable along the whole path, and compares the two martingales it produces through their bracket up to the separation of the two leaves.
\end{abstract}
\maketitle

\section{Introduction}\label{sec:intro}

The critical beta-splitting tree is the case $\beta=-1$ of Aldous's beta-splitting model of phylogenetic trees. Its continuous-time version $\mathrm{CTCS}(n)$ has the following interval-splitting description \cite{AJ,AP}. An interval of $m\ge2$ leaves, labelled $1,\dots,m$ from left to right, is cut between $i$ and $i+1$ at rate
\begin{equation}\label{eq:rate}
\lambda_m(i)=\frac{m}{2i(m-i)},\qquad 1\le i\le m-1,
\end{equation}
so that it waits an exponential time of rate $h_{m-1}=\sum_{i=1}^{m-1}\lambda_m(i)=\sum_{q=1}^{m-1}1/q$ and then splits into the intervals $\{1,\dots,i\}$ and $\{i+1,\dots,m\}$, which evolve independently and are relabelled $1,\dots,i$ and $1,\dots,m-i$. An interval of one leaf does not split. Starting from the interval $\{1,\dots,n\}$, the height $H_j$ of leaf $j$ is the time at which $j$ becomes a singleton, and $D_n=H_J$ for $J$ uniform on $\{1,\dots,n\}$ and independent of the tree. The height $D_n$ is known in detail: $\E D_n=\zeta(2)^{-1}\log n+O(1)$ and $\Var D_n\sim\sigma^2\log n$ with
\[
\sigma^2=\frac{2\zeta(3)}{\zeta(2)^3},
\]
see Aldous and Pittel \cite[Theorem~1.1]{AP}; the constant terms are found by Aldous and Janson \cite{AJ4}, and Kolesnik \cite{Kol} proves a central limit theorem.

Aldous and Janson \cite[Section~3.9]{AJ} ask how the heights of two leaves of one tree are related. They close the interval into a cycle and consider two leaves at distance $r_n$ on it. Concretely, let $J$ be uniform on $\{1,\dots,n\}$, let $J'\in\{1,\dots,n\}$ be congruent to $J+r$ modulo $n$, and put $D^{(1)}_n=H_J$, $D^{(2)}_n=H_{J'}$. Each of the two is distributed as $D_n$, and
\[
\rho(n,r)=\frac{\E[D^{(1)}_nD^{(2)}_n]-(\E D_n)^2}{\Var D_n}
\]
is their correlation coefficient. A Brownian approximation of the two lineages suggests that $\rho(n,r_n)$ should be close to $(\log n-\log r_n)/\log n$, and Aldous and Janson state:

\medskip\noindent\textbf{Open Problem 1} (\cite[Section~3.9]{AJ}). Prove that, for $\log r_n/\log n\to\alpha\in[0,1]$, we have $\rho(n,r_n)\to1-\alpha$.
\medskip

They add that the case $\alpha=1$ is essentially the result of Aldous and Pittel on two uniformly random leaves, cited in \cite{AJ} as Theorem~2.6 of the published version of \cite{AP}. The distance between two leaves of a cycle of length $n$ is at most $n/2$, and we take $1\le r_n\le n/2$ throughout; Remark~\ref{rem:oriented} treats an arbitrary offset.

\begin{theorem}\label{thm:main}
There is an absolute constant $C$ such that for all $n\ge2$ and all integers $1\le r\le n/2$,
\begin{align}
\Bigl|\Cov\bigl(D^{(1)}_n,D^{(2)}_n\bigr)-\sigma^2\log\frac nr\Bigr|&\le C\sqrt{\log n},\label{eq:cov}\\
\bigl|\Var D_n-\sigma^2\log n\bigr|&\le C\sqrt{\log n}.\label{eq:var}
\end{align}
\end{theorem}

\begin{corollary}\label{cor:rho}
Uniformly in $1\le r\le n/2$, $\rho(n,r)=1-\log r/\log n+O((\log n)^{-1/2})$. In particular $\rho(n,r_n)\to1-\alpha$ whenever $\log r_n/\log n\to\alpha\in[0,1]$.
\end{corollary}

\begin{proof}
Since $|\rho|\le1$, it suffices to take $n$ large. By \eqref{eq:var}, $\Var D_n\ge\sigma^2\log n/2$ for $n\ge n_0$. Dividing \eqref{eq:cov} by \eqref{eq:var} gives $\rho(n,r)=(\log(n/r)+O(\sqrt{\log n}))/(\log n+O(\sqrt{\log n}))=1-\log r/\log n+O((\log n)^{-1/2})$, since $0\le\log(n/r)\le\log n$.
\end{proof}

\begin{remark}\label{rem:oriented}
The proof gives, for every $1\le r\le n-1$, the estimate \eqref{eq:cov} with $\log(n/r)$ replaced by $\log(n/\min(r,n-r))$; see the end of Section~\ref{sec:proof}. The leaf at offset $n-1$ is the left neighbour of $J$, so for an offset measured in one direction only the statement of Open Problem~1 needs $r_n\le n/2$, or $\min(r_n,n-r_n)$ in place of $r_n$.
\end{remark}

\begin{remark}
Theorem~\ref{thm:main} also identifies the covariance with the mean time the two leaves spend in a common clade: if $S$ is the time of the cut that separates $J$ from $J'$, then $\E S=\zeta(2)^{-1}\log(n/r)+O(1)$ (Proposition~\ref{prop:sep} and \eqref{eq:ES}), and $\Cov(D^{(1)}_n,D^{(2)}_n)=2\zeta(3)\zeta(2)^{-2}\,\E S+O(\sqrt{\log n})$ by \eqref{eq:covS}. The error $O(\sqrt{\log n})$ comes from the Cauchy--Schwarz inequality in Section~\ref{sec:proof}; the simulations of Section~\ref{sec:numerics} suggest that the true error in \eqref{eq:cov} and \eqref{eq:var} is $O(1)$, and hence that $\rho(n,r)-(1-\log r/\log n)$ is of order $1/\log n$.
\end{remark}

\subsection*{Idea of the proof}
For a leaf at position $j$ in a clade of $m$ leaves, the quantity $F(m,j)=\log(j(m+1-j)/m)$ lies between $0$ and $\log m$ and vanishes at a singleton. Under the dynamics it decreases with drift $\zeta(2)$ and with bracket rate $2\zeta(3)$, up to errors that are small unless the leaf is close to an end of its clade (Lemma~\ref{lem:drift}), and the time the leaf spends close to an end is controlled in $L^2$ (Lemma~\ref{lem:errors}). Thus $\zeta(2)D=F(n,J)+M+O_{L^2}(1)$ for a martingale $M$ (Lemma~\ref{lem:mart}). The heights of the two leaves are then compared through their martingales $M^{(1)},M^{(2)}$: before the two leaves separate their potentials move together, and the difference $M^{(1)}-M^{(2)}$ at the separating cut is bounded in $L^2$ (Lemma~\ref{lem:rect}). Polarization of the brackets shows that the cross bracket accumulated before separation is $2\zeta(3)\E S+O(1)$, and after separation the two martingales have no common jumps. The mean separation time $\E S$ is computed from an exact description of the common clade of the two leaves (Lemma~\ref{lem:pair}, Proposition~\ref{prop:sep}).

\subsection*{Notation}
All logarithms are natural, $h_q=\sum_{t=1}^q1/t$ with $h_0=0$, and $c_1=\zeta(2)$, $c_2=2\zeta(3)$, so that $\sigma^2=c_2/c_1^3$. The letter $C$ denotes an absolute constant whose value may change from line to line. For a random variable $X$, $\|X\|_2=(\E X^2)^{1/2}$.

\section{The positional potential of one leaf}\label{sec:potential}

Consider a tagged leaf at position $j$ in a clade of $m$ leaves, and put
\[
a=j,\qquad b=m+1-j,\qquad k=\min(a,b),\qquad F(m,j)=\log\frac{ab}{m}.
\]
Since $a+b=m+1$ we have $ab-m=(a-1)(b-1)\ge0$, so $0\le F(m,j)\le\log m$, and $F(m,j)=0$ exactly when $k=1$. A cut at $i<a$ removes $l=i$ leaves on the left of the tag and moves it to position $j-l$ in a clade of $m-l$ leaves; a cut at $i\ge a$ removes $l=m-i$ leaves on the right, $1\le l\le b-1$. In both cases the rate of the cut is $\lambda_m(l)=m/(2l(m-l))$, the two families together are all $m-1$ cuts, and the loss $F(\text{before})-F(\text{after})$ is
\begin{equation}\label{eq:loss}
\Delta_{m,a}(l)=\log\frac{a(m-l)}{m(a-l)}\quad(1\le l<a),\qquad \Delta_{m,b}(l)=\log\frac{b(m-l)}{m(b-l)}\quad(1\le l<b),
\end{equation}
which is nonnegative. For $p=1,2$ put
\[
d_p(m,j)=\sum_{l=1}^{a-1}\lambda_m(l)\Delta_{m,a}(l)^p+\sum_{l=1}^{b-1}\lambda_m(l)\Delta_{m,b}(l)^p .
\]

\begin{lemma}\label{lem:drift}
For all $m\ge1$ and $1\le j\le m$, $d_p(m,j)=c_p-e_p(m,j)$ with
\[
0\le e_p(m,j)\le C\,\frac{1+\log^p(k+1)}{k},\qquad p=1,2.
\]
If $k=1$, then $d_1=d_2=0$.
\end{lemma}

\begin{proof}
If $k=1$, say $a=1$, then $b=m$ and every loss $\Delta_{m,m}(l)=\log\frac{m(m-l)}{m(m-l)}$ vanishes; so $d_p=0$ and $e_p=c_p$, which is within the bound. Let $k\ge2$, so $2\le a\le m-1$. For $0<x<a$ put
\[
f_p(x)=\frac{m}{2x(m-x)}\Bigl(\log\frac{a(m-x)}{m(a-x)}\Bigr)^p .
\]
The substitution $t=x(m-a)/(a(m-x))$ maps $(0,a)$ onto $(0,1)$, with $1-t=m(a-x)/(a(m-x))$ and $\dd t/t=m\dd x/(x(m-x))$. Hence
\begin{equation}\label{eq:int}
\int_0^af_p(x)\dd x=\frac12\int_0^1\bigl(-\log(1-t)\bigr)^p\frac{\dd t}{t}=\frac{c_p}2 ,
\end{equation}
because $\int_0^1(-\log(1-t))t^{-1}\dd t=\sum_{q\ge1}q^{-2}=\zeta(2)$ and $\int_0^1\log^2(1-t)\,t^{-1}\dd t=\int_0^1\log^2u\,(1-u)^{-1}\dd u=\sum_{q\ge0}2/(q+1)^3=2\zeta(3)$, the series having nonnegative terms.

The function $f_p$ is increasing on $(0,a)$. Indeed $\log\frac{a(m-x)}{m(a-x)}=\log(1-x/m)-\log(1-x/a)=\sum_{q\ge1}(a^{-q}-m^{-q})x^q/q$ has nonnegative coefficients and no constant term, so its $p$-th power is $x^p$ times a power series with nonnegative coefficients, and $m/(2x(m-x))=(2x)^{-1}\sum_{q\ge0}(x/m)^q$. Comparing $f_p(l)$ with the integral of $f_p$ over $[l-1,l]$ and over $[l,l+1]$ gives
\[
0\le\int_0^af_p-\sum_{l=1}^{a-1}f_p(l)\le\int_{a-1}^af_p .
\]
Write $d'=m-a\ge1$ and $q=\min(a,d')$. For $x=a-u$ with $0<u<1$ we have $\frac{m}{2x(m-x)}=\frac12(\frac1x+\frac1{m-x})\le\frac1a+\frac1{d'}\le\frac2q$ (using $x>a-1\ge a/2$), and $\frac{a(m-x)}{m}\le\min(a,m-x)\le\min(a,d'+1)\le q+1$, so the logarithm in $f_p$ is at most $\log((q+1)/u)$. Therefore
\[
\int_{a-1}^af_p\le\frac2q\int_0^1\log^p\frac{q+1}u\dd u\le C\,\frac{1+\log^p(q+1)}{q}.
\]
The sum over left cuts in $d_p$ is $\sum_{l=1}^{a-1}f_p(l)$, and the same argument with $a$ replaced by $b$ applies to the right cuts. Adding, $e_p=c_p-d_p$ is nonnegative and at most the sum of the two bounds, where $q=\min(a,b-1)$ on the left and $q=\min(b,a-1)$ on the right. Both are at least $k-1\ge k/2$. The function $x\mapsto(1+\log^p(x+1))/x$ is decreasing on $(0,\infty)$ (for $p=2$ the numerator of its derivative is at most $2\ell-1-\ell^2\le0$ with $\ell=\log(x+1)$), so each bound is at most its value at $k/2$, which is at most $2(1+\log^p(k+1))/k$ times the constant.
\end{proof}

Let $(m_t,j_t)$ be the clade size and the position of the tagged leaf at time $t$, started from $(m,j)$, and let $D$ be the time at which $m_t=1$. Write $F_t=F(m_t,j_t)$, $k_t=\min(j_t,m_t+1-j_t)$ and $e_p(t)=e_p(m_t,j_t)$.

\begin{lemma}\label{lem:errors}
For $p=1,2$ and every stopping time $\tau\le D$, $\bigl\|\int_0^\tau e_p(t)\dd t\bigr\|_2\le C$, uniformly in $m$ and $j$.
\end{lemma}

\begin{proof}
Since $e_p\ge0$ it suffices to take $\tau=D$. A cut never increases $k_t$: a left cut replaces $a$ by $a-l$ and keeps $b$, and a right cut does the opposite. For $q\ge1$ let $T_q$ be the time spent with $k_t\in[2^q,2^{q+1})$ and let $T_0$ be the time spent with $k_t=1$ before $D$. While $k_t=k\in[L,2L)$, $L=2^q$, say with $k=a\le b$, each of the $\lfloor k/2\rfloor$ cuts removing $l\in[\lceil k/2\rceil,k-1]$ leaves on the near side has rate $\lambda_m(l)\ge1/(2l)\ge1/(2(k-1))$ and brings $k_t$ to at most $\lfloor k/2\rfloor<L$, after which $k_t$ never returns to $[L,2L)$. Their total rate is at least $\lfloor k/2\rfloor/(2(k-1))\ge1/4$. Hence, whatever happens inside the range, $\Pp(T_q>s)\le e^{-s/4}$, and $\|T_q\|_2\le\sqrt{32}$. While $k_t=1$ and $m_t\ge2$, the cut that isolates the tag has rate $\lambda_m(1)\ge1/2$, so $\|T_0\|_2\le\sqrt8$. On $[2^q,2^{q+1})$, Lemma~\ref{lem:drift} bounds $e_p$ by $C(1+(q+2)^p)2^{-q}$, and $e_p\le c_p$ when $k_t=1$. By Minkowski's inequality,
\[
\Bigl\|\int_0^De_p(t)\dd t\Bigr\|_2\le c_p\|T_0\|_2+\sum_{q\ge1}C(1+(q+2)^p)2^{-q}\|T_q\|_2\le C .\qedhere
\]
\end{proof}

\begin{lemma}\label{lem:mart}
The process $M_t=F_t-F_0+\int_0^{t\wedge D}d_1(m_u,j_u)\dd u$ is a square-integrable martingale with predictable quadratic variation $\int_0^{t\wedge D}d_2(m_u,j_u)\dd u$, and
\begin{equation}\label{eq:mart}
c_1D=F(m,j)+M_D+\int_0^De_1(t)\dd t,\qquad \E M_D^2=c_2\E D-\E\int_0^De_2(t)\dd t .
\end{equation}
In particular $\E D=F(m,j)/c_1+O(1)$ and $\E M_D^2\le C(1+\log m)$.
\end{lemma}

\begin{proof}
The pair $(m_t,j_t)$ is a Markov chain on a finite set, and by \eqref{eq:loss} its generator applied to $(m',j')\mapsto F(m',j')$ equals $-d_1(m',j')$, while the jump rates weighted by the squared jumps of $F$ give $d_2(m',j')$. Hence $M$ is a martingale with the stated predictable quadratic variation. The chain makes at most $m-1$ jumps and every holding rate is $h_{m'-1}\ge1$, so $D$ is dominated by a sum of $m-1$ independent standard exponential variables and has moments of all orders; together with $|F|\le\log m$ this justifies optional stopping at $D$ and the isometry $\E M_D^2=\E\int_0^Dd_2$. Since $F_D=0$ and $d_p=c_p-e_p$, this is \eqref{eq:mart}. Taking expectations and using Lemma~\ref{lem:errors} gives $\E D=F(m,j)/c_1+O(1)$, and then $\E M_D^2\le c_2\E D\le C(1+\log m)$.
\end{proof}

\section{The common clade of two leaves}\label{sec:pair}

Fix $1\le r<n$ and consider two tagged leaves at positions $j$ and $j+r$ of $\{1,\dots,n\}$, where $j$ is uniform on $\{1,\dots,n-r\}$. Let $S$ be the time of the cut that separates them, and for $t<S$ let $Q_t>r$ be the size of the clade containing both, and $j_t$ the position of the left one in it.

\begin{lemma}\label{lem:pair}
Up to time $S$, $(Q_t)$ is a Markov chain that jumps from $m$ to $k$, $r<k<m$, at rate
\begin{equation}\label{eq:arate}
a_r(m,k)=\frac{m(k-r)}{(m-r)k(m-k)},\qquad\text{and}\qquad \frac1{m-k}-a_r(m,k)=\frac{r}{(m-r)k},
\end{equation}
and is killed (the leaves separate) at rate
\begin{equation}\label{eq:kill}
\kappa_r(m)=h_{m-1}-h_{m-r-1}+\frac r{m-r}\,(h_{m-1}-h_r).
\end{equation}
Given $(Q_u)_{u\le t}$ and $t<S$, $j_t$ is uniform on $\{1,\dots,Q_t-r\}$. Given $Q_{S-}=m$ and the past of $Q$, the separating cut is $i=j_{S-}+\theta-1$, and the pair $(j_{S-},\theta)$ has law proportional to $1/(i(m-i))$ on $\{1,\dots,m-r\}\times\{1,\dots,r\}$.
\end{lemma}

\begin{proof}
Suppose $Q_t=m$ and $j_t$ is uniform on $\{1,\dots,m-r\}$. The cut at $i$ keeps both leaves in the left child $\{1,\dots,i\}$ when $j_t\le i-r$, which happens for $(i-r)_+$ values of $j_t$, and keeps both in the right child when $j_t\ge i+1$, for $(m-r-i)_+$ values; in the right child the position becomes $j_t-i$. In either case the retained values of $j_t$ are in bijection with the positions $1,\dots,k-r$ of the new clade of size $k$, so the new position is again uniform. A clade of size $k$ arises from the cut $i=k$ (left child) and from the cut $i=m-k$ (right child), each of rate $m/(2k(m-k))$ and each retaining $k-r$ of the $m-r$ positions; this gives $a_r(m,k)$ (when $k=m-k$ the single cut $i=k$ retains $k-r$ positions on each side, with the same total). The identity in \eqref{eq:arate} is algebra. The holding rate $h_{m-1}$ and the cut rates do not depend on $j_t$, so conditioning on the holding times and on the new size keeps the position uniform, and induction over the jumps gives the uniformity. The killing rate is $h_{m-1}-\sum_{k=r+1}^{m-1}a_r(m,k)$; by \eqref{eq:arate}, $\sum_{k=r+1}^{m-1}a_r(m,k)=h_{m-r-1}-\frac r{m-r}(h_{m-1}-h_r)$, which is \eqref{eq:kill}. Finally the separating cuts are $i=j_t+\theta-1$ with $1\le\theta\le r$; each value of $j_t$ has probability $1/(m-r)$ and each cut has rate $\lambda_m(i)=m/(2i(m-i))$, which gives the last statement.
\end{proof}

\begin{proposition}\label{prop:sep}
Uniformly in $1\le r<n$, $\E S=\zeta(2)^{-1}\log(n/r)+O(1)$.
\end{proposition}

\begin{proof}
Let $g(m)=\log(m/r)$ for $m>r$ and $g=0$ after killing. Its expected rate of decrease at $Q_t=m$ is
\[
\delta_r(m)=\sum_{k=r+1}^{m-1}a_r(m,k)\log\frac mk+\kappa_r(m)\log\frac mr .
\]
Write $a_r(m,k)=\frac1{m-k}-\frac r{(m-r)k}$ and $\kappa_r(m)=\sum_{k=1}^r\frac1{m-k}+\frac r{m-r}\sum_{k=r+1}^{m-1}\frac1k$. Collecting terms,
\begin{equation}\label{eq:drift2}
\delta_r(m)=\delta_0(m)-\sum_{k=1}^r\frac{\log(r/k)}{m-k}+\frac r{m-r}\sum_{k=r+1}^{m-1}\frac{\log(k/r)}k,\qquad \delta_0(m)=\sum_{k=1}^{m-1}\frac{\log(m/k)}{m-k}.
\end{equation}
The function $\phi(x)=-\log x/(1-x)$ is decreasing on $(0,1)$, since $\phi'(x)$ has the sign of $1-1/x-\log x\le0$; it has $\phi(1-)=1$ and $\int_0^1\phi=\sum_{q\ge1}q^{-2}=\zeta(2)$. As $\delta_0(m)=\frac1m\sum_{k=1}^{m-1}\phi(k/m)$, comparing with the integral over the intervals to the left and to the right of each grid point gives
\begin{equation}\label{eq:d0}
0\le\zeta(2)-\delta_0(m)\le\int_0^{1/m}\phi\le\frac{2(1+\log m)}m .
\end{equation}
For $m\ge2r$, $\sum_{k\le r}\log(r/k)=r\log r-\log r!\le r$ and $m-k\ge m/2$ bound the first correction by $2r/m$, and $\log(k/r)\le\log(m/r)$, $\sum_{k=r+1}^{m-1}1/k\le\log(m/r)$ bound the second by $(2r/m)\log^2(m/r)$. Hence
\begin{equation}\label{eq:drift-bound}
|\delta_r(m)-\zeta(2)|\le C\Bigl[\frac{1+\log m}m+\frac rm\Bigl(1+\log^2\frac mr\Bigr)\Bigr],\qquad m\ge2r.
\end{equation}
For $r<m<2r$, $\delta_r(m)\ge0$ because every jump decreases $g$, and the second correction is at most $\frac1{m-r}\sum_{k=r+1}^{m-1}\frac{k-r}k<1$, using $\log(k/r)\le(k-r)/r$; so $0\le\delta_r(m)\le\zeta(2)+1$. Together with \eqref{eq:drift-bound}, $|\delta_r-\zeta(2)|\le C$ for $r<m<4r$.

The size $Q_t$ only decreases. While $r<Q_t<4r$ the killing rate is at least $\sum_{k=1}^r1/(m-k)\ge r/(m-1)>1/4$, so the time spent there has mean at most $4$. For $q\ge0$ and $L=4r2^q$, while $Q_t=m\in[L,2L)$ every target $k\in[\lceil L/2\rceil,L-1]$ satisfies $k\ge2r$, so $a_r(m,k)\ge(k-r)/(k(m-k))\ge1/(2(m-k))>1/(4L)$; there are at least $L/2$ such targets, the exit rate from $[L,2L)$ is at least $1/8$, and the mean time spent there is at most $8$. On $[L,2L)$ the right side of \eqref{eq:drift-bound} is at most $C[(1+\log r+q)/(r2^q)+(1+q^2)2^{-q}]$, which is summable in $q$ with a sum bounded uniformly in $r$. Therefore $\E\int_0^S|\delta_r(Q_t)-\zeta(2)|\dd t\le C$. Since $Q$ has finitely many states and $\E S<\infty$, Dynkin's formula gives $\log(n/r)=g(n)=\E\int_0^S\delta_r(Q_t)\dd t=\zeta(2)\E S+O(1)$.
\end{proof}

\section{The separating cut}\label{sec:cut}

Let $r,s\ge1$, $m=s+r$, and on the rectangle $R=\{1,\dots,s\}\times\{1,\dots,r\}$ consider the probability $w$ proportional to $1/(i(m-i))$, where $i=j+\theta-1$ for $(j,\theta)\in R$. Put $b=r+1-\theta$, $h=s+1-j$ and
\[
G(j,\theta)=\log\frac{j\,\theta\,(m-i)}{i\,b\,h}.
\]
The left leaf of Section~\ref{sec:pair}, at position $j$ of a clade of $m$ leaves, is sent by the cut $i=j+\theta-1$ to position $j$ of a clade of $i$ leaves, where $j+\theta-1=i$ means that its two ranks are $j$ and $\theta$; the right leaf, at $j+r$, goes to position $j+r-i=b$ of a clade of $m-i$ leaves, where its rank from the right is $m-i+1-b=h$. So $G(j,\theta)$ is the difference of the potentials of the two leaves just after the separating cut.

\begin{lemma}\label{lem:rect}
$\E_wG^2\le C$, uniformly in $r,s\ge1$.
\end{lemma}

\begin{proof}
The weight and $G$ are unchanged when $(j,s)$ and $(\theta,r)$ are exchanged, so we may assume $s\ge r$. The reflection $(j,\theta)\mapsto(h,b)$ maps $R$ onto itself, sends $i$ to $m-i$, preserves $w$ and changes the sign of $G$. Since $1/(i(m-i))=(1/i+1/(m-i))/m$, the total mass is $\frac2m Z$ with $Z=\sum_R1/i$, and for every reflection-invariant $Q$, $\E_wQ=Z^{-1}\sum_RQ/i$.

Let $L=\log(1+s/r)\ge\log2$ and $H(\theta)=\sum_{j=1}^s1/(j+\theta-1)$. Comparing with integrals,
\[
L\le\log\frac{s+\theta}\theta\le H(\theta)\le1+\log\frac{s+r}\theta=1+L+\log\frac r\theta,\qquad Z=\sum_{\theta=1}^rH(\theta)\ge rL .
\]
Write $G=P-A+B$ with $P=\log(\theta/b)$, $A=\log(i/j)$, $B=\log((m-i)/h)$. The reflection changes the sign of $P$ and exchanges $A$ and $B$, so $\E_wB^2=\E_wA^2$, and $\E_wG^2\le3\E_wP^2+6\E_wA^2$.

For a decreasing $\psi\ge0$ on $(0,1]$, $\frac1r\sum_{\theta=1}^r\psi(\theta/r)\le\int_0^1\psi$; with $\psi=\log^p(1/x)$ this gives $\frac1r\sum_\theta\log^p(r/\theta)\le p!$, and the same for $\log^p(r/b)$, since $b$ also runs over $1,\dots,r$. As $|P|\le\log(r/\theta)+\log(r/b)$ and $xy^2\le(x^3+2y^3)/3$ for $x,y\ge0$,
\[
\sum_{\theta=1}^rH(\theta)P^2\le\sum_{\theta=1}^r\Bigl(1+L+\log\frac r\theta\Bigr)\cdot2\Bigl(\log^2\frac r\theta+\log^2\frac rb\Bigr)\le Cr(1+L),
\]
so $\E_wP^2=Z^{-1}\sum_\theta H(\theta)P^2\le C(1+L)/L\le C$.

Next, $0\le A=\log(1+(\theta-1)/j)\le\log(1+\theta/j)$ and $\E_wA^2=(2Z)^{-1}\sum_RA^2(1/i+1/(m-i))$. For the first part, $i\ge(j+\theta)/2$, and $x\mapsto\log^2(1+\theta/x)/(x+\theta)$ is decreasing, so
\[
\sum_R\frac{A^2}i\le2\sum_{\theta=1}^r\int_0^\infty\frac{\log^2(1+\theta/x)}{x+\theta}\dd x=2r\int_0^\infty\frac{\log^2(1+1/u)}{1+u}\dd u\le Cr ,
\]
the last integral being finite (it behaves like $\log^2u$ near $0$ and like $u^{-3}$ at infinity). For the second part, $m-i=s-j+b$. If $j\le s/2$ then $m-i\ge s/2$ and $\sum_{j\ge1}\log^2(1+\theta/j)\le\theta\int_0^\infty\log^2(1+1/u)\dd u\le C\theta$, which contributes at most $C\sum_\theta\theta/s\le Cr^2/s$. If $j>s/2$ then $A\le\theta/j\le2\theta/s$ and $\sum_{j>s/2}1/(s-j+b)\le1+\log(1+s/b)\le1+L+\log(r/b)$, which contributes at most $\sum_\theta\frac{4\theta^2}{s^2}(1+L+\log\frac rb)\le C\frac{r^3}{s^2}(1+L)$. With $r\le s$ both contributions are at most $Cr(1+L)$, and $\E_wA^2\le C(1+L)/L\le C$.
\end{proof}

\section{Proof of Theorem~\ref{thm:main}}\label{sec:proof}

Run $\mathrm{CTCS}(n)$ with the two leaves $J$ and $J'$ of Section~\ref{sec:intro}, and let $\mathcal F_t$ be generated by $J$, $J'$ and the cuts made up to time $t$. For $\nu=1,2$ let $(m^{(\nu)}_t,j^{(\nu)}_t)$ be the clade size and the position of the $\nu$-th leaf, $D_\nu$ its height, $F^{(\nu)}_t$ its potential, $e^{(\nu)}_p(t)$ its error terms and
\[
M^{(\nu)}_t=F^{(\nu)}_t-F^{(\nu)}_0+\int_0^{t\wedge D_\nu}d_1\bigl(m^{(\nu)}_u,j^{(\nu)}_u\bigr)\dd u .
\]
The clade containing a given leaf is cut at the rates \eqref{eq:rate} whatever the rest of the tree does, so the proof of Lemma~\ref{lem:mart} applies in the filtration $(\mathcal F_t)$: each $M^{(\nu)}$ is an $(\mathcal F_t)$-martingale with predictable quadratic variation $\int d_2(m^{(\nu)}_u,j^{(\nu)}_u)\dd u$, and \eqref{eq:mart} holds for it. Lemma~\ref{lem:errors} holds for each leaf separately.

\emph{The initial potentials.} With $J$ uniform, $F(n,J)-\log n=\log(J/n)+\log((n+1-J)/n)$, and $\E\log^2(n/J)\le\int_0^1\log^2x\dd x=2$ by the comparison used in Lemma~\ref{lem:rect}; the same holds for $n+1-J$ and for $J'$. Hence $\|F^{(\nu)}_0-\log n\|_2\le C$.

\emph{Reduction to the separation time.} Let $S$ be the time of the cut that separates the two leaves. If $J\le n-r$, an event of probability $(n-r)/n$, the leaves are $J<J'=J+r$ and, given this event, $J$ is uniform on $\{1,\dots,n-r\}$; otherwise they are $J'=J+r-n<J$, at linear distance $n-r$, and $J'$ is uniform on $\{1,\dots,r\}$. In both cases Section~\ref{sec:pair} applies to the ordered pair, so by Proposition~\ref{prop:sep}
\begin{equation}\label{eq:ES}
\E S=\frac1{\zeta(2)}\Bigl[\frac{n-r}n\log\frac nr+\frac rn\log\frac n{n-r}\Bigr]+O(1)=\frac{\log(n/r)}{\zeta(2)}+O(1),
\end{equation}
using $r\le n/2$, $\frac rn\log\frac nr\le1$ and $\log\frac n{n-r}\le\log2$.

\emph{The difference at separation.} From here on $\nu=1$ is the left leaf of the ordered pair and $\nu=2$ the right one. The bound $\|F^{(\nu)}_0-\log n\|_2\le C$ still holds, since $(F^{(1)}_0-\log n)^2+(F^{(2)}_0-\log n)^2$ is the same sum for the original labels. Let $N_t=M^{(1)}_{t\wedge S}-M^{(2)}_{t\wedge S}$. Before $S$ both leaves have $d_1=c_1-e_1$, so
\[
N_S=\bigl(F^{(1)}_S-F^{(2)}_S\bigr)-\bigl(F^{(1)}_0-F^{(2)}_0\bigr)-\int_0^Se^{(1)}_1(t)\dd t+\int_0^Se^{(2)}_1(t)\dd t .
\]
By Lemma~\ref{lem:pair}, given the separation from a clade of size $m$, the cut and the position of the left leaf have the law $w$ of Section~\ref{sec:cut} with $s=m-r$ (or $m-(n-r)$), and $F^{(1)}_S-F^{(2)}_S=G$; hence $\E(F^{(1)}_S-F^{(2)}_S)^2\le C$ by Lemma~\ref{lem:rect}, after averaging over $m$. With the initial bound and Lemma~\ref{lem:errors},
\begin{equation}\label{eq:NS}
\E N_S^2\le C .
\end{equation}

\emph{Polarization.} While the two leaves share a clade, every cut of that clade moves both potentials; let $\beta_t$ be the sum over these cuts of the rate times the product of the two losses, so that the predictable covariation of $M^{(1)}$ and $M^{(2)}$ is $\int_0^{t\wedge S}\beta_u\dd u$. After $S$ the two leaves lie in disjoint clades, which are cut by independent exponential clocks, so $M^{(1)}$ and $M^{(2)}$ have no common jumps and their covariation does not grow. The martingale $N$ is square integrable, and $\E N_S^2$ equals the expected predictable quadratic variation of $N$ up to $S$:
\[
\E N_S^2=\E\int_0^S\Bigl(d_2\bigl(m^{(1)}_t,j^{(1)}_t\bigr)+d_2\bigl(m^{(2)}_t,j^{(2)}_t\bigr)-2\beta_t\Bigr)\dd t .
\]
Since $d_2=c_2-e_2$, Lemma~\ref{lem:errors} and \eqref{eq:NS} give $\E\int_0^S\beta_t\dd t=c_2\E S+O(1)$. The product $M^{(1)}M^{(2)}$ minus their predictable covariation is a martingale, uniformly integrable because the chain is finite and both heights have all moments, so
\begin{equation}\label{eq:EMM}
\E\bigl[M^{(1)}_{D_1}M^{(2)}_{D_2}\bigr]=\E\int_0^S\beta_t\dd t=c_2\E S+O(1).
\end{equation}

\emph{The covariance.} By \eqref{eq:mart}, $D_\nu=c_1^{-1}\log n+c_1^{-1}M^{(\nu)}_{D_\nu}+R_\nu$ with $R_\nu=c_1^{-1}\bigl(F^{(\nu)}_0-\log n+\int_0^{D_\nu}e^{(\nu)}_1\bigr)$, so $\|R_\nu\|_2\le C$, and $\|M^{(\nu)}_{D_\nu}\|_2\le C\sqrt{\log n}$ by Lemma~\ref{lem:mart}. As $\E M^{(\nu)}_{D_\nu}=0$, the Cauchy--Schwarz inequality and \eqref{eq:EMM} give
\begin{equation}\label{eq:covS}
\Cov(D_1,D_2)=c_1^{-2}\E\bigl[M^{(1)}_{D_1}M^{(2)}_{D_2}\bigr]+O\bigl(\sqrt{\log n}\bigr)=\frac{c_2}{c_1^2}\,\E S+O\bigl(\sqrt{\log n}\bigr),
\end{equation}
and \eqref{eq:ES} turns this into \eqref{eq:cov}, as $c_2/c_1^3=\sigma^2$. This is the covariance of the ordered pair, and it equals that of $D^{(1)}_n$ and $D^{(2)}_n$. Indeed $D_1D_2$ is the product of the two original heights, so it remains to see that $\E D_1=\E D_2$, both then being $\E D_n$. Since $\lambda_m(i)=\lambda_m(m-i)$, the reflection $i\mapsto n+1-i$ of the leaf labels preserves the law of the tree; it maps a pair at linear distance $g$ with left leaf $a$ to the pair at linear distance $g$ with left leaf $n+1-a-g$, which is again uniform on $\{1,\dots,n-g\}$, and it exchanges the roles of the left and the right leaf. Hence, given either of the two events above, $(D_1,D_2)$ and $(D_2,D_1)$ have the same law. In the same way, for one leaf, $\Var D_n=c_1^{-2}\E M_D^2+O(\sqrt{\log n})$, and $\E M_D^2=c_2\E D_n+O(1)=c_2c_1^{-1}\log n+O(1)$ by \eqref{eq:mart} and Lemmas~\ref{lem:errors} and~\ref{lem:mart}. This is \eqref{eq:var}.

For an offset $n/2<r\le n-1$ the argument is unchanged up to \eqref{eq:ES}, whose bracket now differs from $\log(n/(n-r))$ by $O(1)$; this gives Remark~\ref{rem:oriented}. \qed

\section{Numerical checks}\label{sec:numerics}

The script \texttt{code/checks.py} in the repository below checks the finite statements of the paper and reports 22 passed checks; its output is \texttt{code/checks\_output.txt}. Exact rational arithmetic, by enumeration of all starting positions and cuts for $1\le r<m\le26$, confirms the rates \eqref{eq:arate} and \eqref{eq:kill}, the uniformity of the position in Lemma~\ref{lem:pair}, and the law of the separating cut. In double precision it confirms Lemma~\ref{lem:drift} for all $1\le j\le m\le400$ (the largest values of $e_pk/(1+\log^p(k+1))$ are $1.166$ for $p=1$ and $2.130$ for $p=2$), the identity \eqref{eq:drift2} for $r<m\le600$, the bounds \eqref{eq:d0} and \eqref{eq:drift-bound}, and the rates used in the proofs of Lemma~\ref{lem:errors} and Proposition~\ref{prop:sep}. The exact recursion for the killed chain gives $\E S-\zeta(2)^{-1}\log(n/r)\in[-0.1374,0.1945]$ for $n\in\{2^9,2^{11},2^{13}\}$ and $r\in\{1,4,\dots,4^6\}$ with $r<n$, and $\E D_n-\zeta(2)^{-1}\log n=0.79512$ at $n=2^{13}$, close to the constant $\zeta(3)/\zeta(2)^2+\gamma/\zeta(2)=0.79516$ of \cite{AJ4}. On a grid of $r,s\le3000$ the moment $\E_wG^2$ of Lemma~\ref{lem:rect} is at most $4.64$, while $\E_w\log^2\theta$ grows with $r$.

A direct simulation of the two lineages, which uses none of the potentials, with $40\,000$ trees per row, gives:
\begin{center}
\begin{tabular}{rrcccc}
\toprule
$n$ & $r$ & $\Cov-\sigma^2\log(n/r)$ & $\Var-\sigma^2\log n$ & $\rho(n,r)$ & $1-\alpha$\\
\midrule
$2^{20}$ & $2^{5}$ & $-0.78$ & $0.92$ & $0.575$ & $0.75$\\
$2^{20}$ & $2^{10}$ & $-0.67$ & $0.97$ & $0.364$ & $0.50$\\
$2^{20}$ & $2^{15}$ & $-0.73$ & $0.97$ & $0.135$ & $0.25$\\
$2^{40}$ & $2^{10}$ & $-0.74$ & $0.90$ & $0.661$ & $0.75$\\
$2^{40}$ & $2^{20}$ & $-0.68$ & $1.00$ & $0.426$ & $0.50$\\
$2^{40}$ & $2^{30}$ & $-0.85$ & $0.97$ & $0.181$ & $0.25$\\
\bottomrule
\end{tabular}
\end{center}
The offsets in the third and fourth columns show no growth between $n=2^{20}$ and $n=2^{40}$. This is consistent with an $O(1)$ error, which would make $\rho(n,r)$ approach $1-\alpha$ at rate $1/\log n$; Theorem~\ref{thm:main} proves the rate $(\log n)^{-1/2}$. At $n=2^{20}$ the offsets $1$ and $n-1$ give $\rho=0.890$ and $0.893$.

The Lean~4 file \texttt{lean/PairIdentities.lean} (with Mathlib) checks the algebraic identities behind \eqref{eq:loss}, \eqref{eq:int} and \eqref{eq:arate}, the decomposition of the weight in Lemma~\ref{lem:rect}, the ranks after the separating cut, the counting inequalities used for the exit rates, the inequality $xy^2\le(x^3+2y^3)/3$ and the polarization identity.

\section*{Acknowledgements and funding}
The research presented in this paper was supported by the European Research
Council (ERC) under the European Union's Horizon Europe research and
innovation programme (grant agreement No.~101041711), by the Simons
Foundation as part of the Collaboration on the Mathematical and Scientific
Foundations of Deep Learning, by Heights Labs, by Convex Nexus Capital, by
the Israel Science Foundation (grant number 2258/19), and by the Israel
Science Foundation (ISF Grant 4101/25). Computational resources were
provided in part by Google Cloud through its education credits program
(credit 530550103).

\section*{Declaration of competing interest}
The author holds equity in Heights Labs and in Convex Nexus Capital and
has a management role at Convex Nexus Capital. These entities are included
in the funding declaration above.

\section*{Declaration of generative AI and AI-assisted technologies}
The paper comes out of Math Brain, an ongoing project on the use of AI systems for research mathematics. On 11 October 2026 an agent running in OpenAI's Codex (GPT-6.1 Sol) obtained Theorem~\ref{thm:main} and its proof; another Codex agent and a Claude agent (Anthropic) checked the proof, the latter with an independent simulation. On 14 and 15 October 2026 Claude Opus~5.5, working in Claude Code, checked the proof line by line, wrote the computer checks, the Lean file and the paper in its present form. Two further Codex agents (GPT-6.1 Sol) read the proofs and repeated the computations independently. The author directed the work and is responsible for the content of the paper and for decisions concerning its dissemination.

\section*{Data availability}
No data were used. The code and the Lean file are available at\\ \url{https://github.com/yspennstate/beta-splitting-leaf-correlation}.

\begin{thebibliography}{9}

\bibitem{AJ}
D.~J. Aldous and S.~Janson, The critical beta-splitting random tree II: overview and open problems, arXiv:2303.02529v3 (2025).

\bibitem{AJ4}
D.~J. Aldous and S.~Janson, The critical beta-splitting random tree IV: Mellin analysis of leaf height, \emph{Electron. J. Probab.} \textbf{30} (2025), Paper No.~69, 39~pp.

\bibitem{AP}
D.~J. Aldous and B.~Pittel, The critical beta-splitting random tree I: heights and related results, \emph{Ann. Appl. Probab.} \textbf{35} (2025), 158--195. Theorem~1.1 is cited in the numbering of arXiv:2302.05066v4.

\bibitem{Kol}
B.~Kolesnik, Critical beta-splitting, via contraction, \emph{Electron. Commun. Probab.} \textbf{30} (2025), Paper No.~10, 14~pp.

\end{thebibliography}
\end{document}
